%!TEX root = ../../main.tex
%% ===============================================================
%% 2.3 이진 분류의 확률적 정식화
%% 파일: chapters/02_background/03_classification.tex
%% 상위: chapters/02_background/chapter.tex
%% 정의 라벨: sec:bg-classification, eq:bce
%% 참조 라벨: -
%% 포함 객체: -
%% 인용: vapnik1998statistical
%% 상태: 초안 (docs/OUTLINE.md 참고)
%% ===============================================================

\section{이진 분류의 확률적 정식화}
\label{sec:bg-classification}

\begin{definition}[이진 분류 문제]
입력 공간 $\mathcal{X}$와 라벨 $y \in \{0,1\}$에 대한 결합 분포 $P(X, Y)$가 있을 때, 점수 함수 $f: \mathcal{X} \to \R$를 학습하여 $P(Y=1 \mid X=x)$를 근사하는 문제. 확률 출력은 $p(x) = \sigma(f(x))$, $\sigma(z) = 1/(1+e^{-z})$로 얻는다.
\end{definition}

학습은 경험적 위험 최소화(empirical risk minimization) \cite{vapnik1998statistical}의 틀에서 이루어지며, 본 논문은 로그 손실(이진 교차 엔트로피)을 기본 목적함수로 사용한다:
\begin{equation}
  \mathcal{L}_{\mathrm{BCE}}(p, y) = -\big[y \log p + (1-y)\log(1-p)\big].
  \label{eq:bce}
\end{equation}
